\documentclass[11pt,a4paper]{article}

\usepackage[utf8]{inputenc}
\usepackage[T1]{fontenc}
\usepackage[spanish,es-nodecimaldot]{babel}
\usepackage{amsmath,amssymb}
\usepackage{geometry}
\geometry{margin=2.3cm}
\usepackage{xcolor}
\usepackage{booktabs}
\usepackage{enumitem}
\usepackage{mdframed}
\usepackage[hidelinks]{hyperref}
\usepackage{parskip}

\definecolor{azul}{RGB}{0,80,150}
\definecolor{verde}{RGB}{18,122,18}
\definecolor{rojo}{RGB}{192,0,0}
\definecolor{grispanel}{gray}{0.95}

\newcommand{\sinc}{\operatorname{sinc}}
\newcommand{\Fou}{\mathcal{F}}

% Caja para "lo que digo en voz alta"
\newmdenv[linecolor=azul,linewidth=1pt,backgroundcolor=grispanel,
  roundcorner=4pt,innertopmargin=6pt,innerbottommargin=6pt]{decir}
% Caja para preguntas del profesor
\newmdenv[linecolor=rojo,linewidth=0.8pt,topline=false,bottomline=false,
  rightline=false,innerleftmargin=10pt,leftmargin=2pt]{pregunta}

\title{\textbf{Defensa de los simuladores de difracción}\\[4pt]
\large Cómo implementé los cálculos de Fraunhofer y de Fresnel,\\
figura por figura, y cómo lo explico}
\author{Isaac David --- Óptica, Taller 4 y Parcial 4}
\date{\today}

\begin{document}
\maketitle

\begin{decir}
\textbf{Cómo usar este documento.} No es la teoría formal (esa está en
\texttt{Documentacion\_Fisica.pdf}). Esto es mi \emph{guión de defensa}: lo que
respondo cuando el profesor pregunta ``¿qué hiciste para los cálculos?''. Las
cajas azules son lo que digo en voz alta; las secciones marcadas con
``\textcolor{rojo}{P:}'' son preguntas que probablemente me hará y mi respuesta.
\end{decir}

\tableofcontents
\bigskip

%======================================================================
\section{La respuesta de una sola frase}
%======================================================================

Si el profesor solo me deja decir una cosa, digo esto:

\begin{decir}
Todo el proyecto se apoya en un único resultado: \textbf{el patrón de difracción
es el módulo al cuadrado de la transformada de Fourier de la abertura}. En campo
lejano (Fraunhofer) es la transformada \emph{pura}; en campo cercano (Fresnel) es
la misma transformada pero con la abertura multiplicada por una fase cuadrática
(un ``chirp''). Entonces mi trabajo, para cada figura, se reduce a tres cosas:
\emph{(1)} escribir la función de transmisión $t(x,y)$ de esa abertura,
\emph{(2)} transformarla ---analíticamente si tiene fórmula cerrada, o con la FFT
si no la tiene---, y \emph{(3)} tomar el módulo al cuadrado. La física está toda
en el paso de la transformada; lo demás es álgebra y graficación.
\end{decir}

La única diferencia entre los tres códigos es \emph{cómo} calculo esa transformada:

\begin{center}
\begin{tabular}{@{}lll@{}}
\toprule
\textbf{Código} & \textbf{Qué calcula} & \textbf{Cómo} \\
\midrule
20 & Fraunhofer & Fórmula cerrada de la TF (sinc, Bessel) \\
21 & Fraunhofer & TF numérica: \texttt{fft2} de la máscara \\
22 & Fresnel$\to$Fraunhofer & TF con chirp (\texttt{fresnel\_propagate}), vs.\ $N_F$ \\
\bottomrule
\end{tabular}
\end{center}

%======================================================================
\section{El método general: la receta que aplico a cualquier figura}
%======================================================================

Quiero que quede claro que \textbf{no memoricé una fórmula por figura}: apliqué el
mismo procedimiento cada vez. Esta es la receta:

\begin{enumerate}[leftmargin=1.5em]
\item \textbf{Escribo la abertura} como $t(x,y)$: vale 1 donde pasa la luz y 0
donde el material la bloquea (aberturas binarias). El campo justo detrás es
$U_0(x,y)=t(x,y)$.
\item \textbf{Elijo el régimen} con el número de Fresnel $N_F=D^2/\lambda z$
(sección~\ref{sec:regimen}). $N_F\le0.5\Rightarrow$ Fraunhofer; si no, Fresnel.
\item \textbf{Transformo:}
   \begin{itemize}[leftmargin=1.3em]
   \item Fraunhofer con fórmula: uso los ``bloques'' de la
   sección~\ref{sec:bloques} (rect$\to$sinc, círculo$\to$Airy) y los combino.
   \item Fraunhofer sin fórmula (corazón, triángulo\ldots): FFT de la máscara.
   \item Fresnel: multiplico $t$ por el chirp $e^{i\pi(x^2+y^2)/\lambda z}$ y hago
   un solo FFT (sección~\ref{sec:fresnel}).
   \end{itemize}
\item \textbf{Tomo $|\cdot|^2$} y normalizo.
\item \textbf{Verifico con un caso límite conocido} antes de darlo por bueno
(sección~\ref{sec:validacion}). Este paso es el que convierte ``creo que está
bien'' en ``sé que está bien''.
\end{enumerate}

\begin{decir}
El punto que quiero recalcar: la implementación matemática es \emph{la misma
maquinaria} para las nueve figuras. Lo único que cambia entre una figura y otra es
la expresión de $t(x,y)$ que le meto a la transformada. Por eso puedo defender
cualquiera: todas salen del mismo molde.
\end{decir}

%======================================================================
\section{Fraunhofer: de dónde sale que \texorpdfstring{$I=|\Fou\{t\}|^2$}{I=|FT|2}}
%======================================================================

Esto lo puedo derivar en el tablero si me lo piden. Parto de la integral de
Fresnel (propagación paraxial una distancia $z$):
\begin{equation}
U(x',y') = \frac{e^{ikz}}{i\lambda z}
\iint t(x,y)\,
e^{\frac{i\pi}{\lambda z}\left[(x'-x)^2+(y'-y)^2\right]}\,dx\,dy .
\label{eq:fresnel}
\end{equation}
Desarrollo el cuadrado $(x'-x)^2 = x'^2 - 2x'x + x^2$ y saco de la integral todo lo
que no depende de $(x,y)$:
\begin{equation}
U(x',y') = \frac{e^{ikz}}{i\lambda z}\,
e^{\frac{i\pi}{\lambda z}(x'^2+y'^2)}
\iint \underbrace{t(x,y)\,e^{\frac{i\pi}{\lambda z}(x^2+y^2)}}_{\text{chirp de entrada}}\;
e^{-\frac{i2\pi}{\lambda z}(x'x+y'y)}\,dx\,dy .
\label{eq:fresnelfft}
\end{equation}
Esa última integral \textbf{es una transformada de Fourier} evaluada en las
frecuencias
\begin{equation}
f_x=\frac{x'}{\lambda z},\qquad f_y=\frac{y'}{\lambda z}.
\label{eq:mapeo}
\end{equation}

\textbf{El paso a Fraunhofer:} cuando $z$ es grande, la fase del chirp de entrada
$e^{i\pi(x^2+y^2)/\lambda z}$ es despreciable sobre toda la abertura (su valor
$\to 1$). Entonces desaparece y queda
\begin{equation}
\boxed{\,U(x',y')\propto \Fou\{t(x,y)\}\big(f_x,f_y\big),\qquad I=|U|^2\, .}
\end{equation}

\begin{decir}
Fraunhofer no es una fórmula distinta de Fresnel: es Fresnel cuando el chirp de
entrada vale 1. Por eso en el código el mismo motor sirve para los dos ---la
diferencia es si dejo o no el factor de fase cuadrático. Y el mapeo
$f_x=x'/\lambda z$ es el que convierte ``frecuencia espacial'' en ``posición en la
pantalla''; es la línea \texttt{fx = X/(lam*z)} que aparece en cada función.
\end{decir}

%======================================================================
\section{Los cuatro bloques matemáticos de Fraunhofer}
\label{sec:bloques}
%======================================================================

Todas las figuras analíticas (Código 20) las armé combinando cuatro piezas. Si
domino estas cuatro, domino las nueve figuras.

\paragraph{Bloque 1 --- Rectángulo $\to$ sinc.}
La TF de un rectángulo de lados $a\times b$ es
\[
\Fou\{\text{rect}_{a\times b}\}=ab\,\sinc(a f_x)\,\sinc(b f_y),
\]
con la sinc \emph{normalizada} $\sinc(u)=\sin(\pi u)/(\pi u)$. Primer cero (primer
mínimo de difracción) en $x'=\lambda z/a$. En código:
\texttt{a*b*np.sinc(a*fx)*np.sinc(b*fy)}.

\paragraph{Bloque 2 --- Círculo $\to$ disco de Airy.}
Por simetría circular, la TF de un disco de radio $R$ es una Bessel de orden 1:
\[
\Fou\{\text{circ}_R\}=\pi R^2\,\frac{2 J_1(2\pi R\rho)}{2\pi R\rho},
\qquad \rho=\sqrt{f_x^2+f_y^2},
\]
con límite $\pi R^2$ (el área) en $\rho\to0$. Primer anillo oscuro donde
$J_1$ se anula: $2\pi R\rho=3.8317\Rightarrow x'=1.22\,\lambda z/(2R)$. En código
uso \texttt{scipy.special.j1} y trato el $0/0$ del origen con \texttt{np.where}.

\paragraph{Bloque 3 --- Desplazamiento $\to$ interferencia.}
Si tengo dos sub-aberturas separadas una distancia $D$, el teorema de
desplazamiento de Fourier dice que trasladar $\pm D/2$ multiplica la TF por
$e^{\mp i\pi D f_x}$. Al sumar las dos amplitudes y tomar $|\cdot|^2$ aparece el
término cruzado:
\begin{equation}
I = A_1^2 + A_2^2 + 2A_1 A_2\cos(2\pi D f_x).
\label{eq:interf}
\end{equation}
Ese coseno son las \textbf{franjas de interferencia} de período angular
$\lambda/D$ ---la misma física de la doble rendija de Young, pero con
sub-aberturas de forma arbitraria. En código: \texttt{2*A1*A2*np.cos(2*np.pi*D*fx)}.

\paragraph{Bloque 4 --- Linealidad $\to$ uniones e inclusión-exclusión.}
La TF es lineal, así que la transformada de una unión de formas es la suma de sus
transformadas ---restando los solapes para no contarlos dos veces
($\mathbb{1}_{A\cup B}=\mathbb{1}_A+\mathbb{1}_B-\mathbb{1}_{A\cap B}$). Con esto
armo el marco (exterior $-$ interior) y la cruz (barra $+$ barra $-$ cuadrado
central).

%======================================================================
\section{Fraunhofer figura por figura (Código 20)}
%======================================================================

Para cada figura muestro la misma cadena: \emph{abertura $\to$ expresión
$\to$ cómo la validé}. Todo sale de combinar los cuatro bloques.

\subsection{Marco + círculo (Punto 1 del Parcial)}
\textbf{Abertura:} marco rectangular en $x=-D/2$ y círculo en $x=+D/2$.
\textbf{Expresión} (bloques 1, 2, 3):
\[
A_\text{marco}=ab\,\sinc(af_x)\sinc(bf_y)-cd\,\sinc(cf_x)\sinc(df_y),\quad
A_\text{c}=\pi R^2\tfrac{2J_1(2\pi R\rho)}{2\pi R\rho},
\]
\[
I=A_\text{marco}^2+A_\text{c}^2+2A_\text{marco}A_\text{c}\cos(2\pi D f_x).
\]
\textbf{Validación:} normalizo a $I(0,0)=1$ con $A_0=ab-cd+\pi R^2$; compruebo que
el paso de las franjas verticales es $\lambda z/D$. (Función \texttt{intensidad}.)

\subsection{Rectángulo rotado (Punto 3)}
\textbf{Idea clave:} rotar la abertura un ángulo $\theta$ equivale a rotar el plano
de frecuencias el mismo ángulo (propiedad de la TF: $\Fou\{f(R_\theta r)\}=
F(R_\theta k)$). \textbf{Expresión:}
\[
f_a=f_x\cos\theta+f_y\sin\theta,\;\; f_b=-f_x\sin\theta+f_y\cos\theta,
\qquad I=\sinc^2(af_a)\,\sinc^2(bf_b).
\]
\textbf{Validación:} en $\theta=0$ recupera el rectángulo alineado; el patrón
sinc$^2$ simplemente gira solidario con la abertura.

\subsection{Cruz (Punto 5)}
\textbf{Abertura:} barra horizontal $L\times a$ $\cup$ barra vertical $a\times L$.
Por inclusión-exclusión (bloque 4):
\[
A_\text{cruz}=La\,\sinc(Lf_x)\sinc(af_y)+aL\,\sinc(af_x)\sinc(Lf_y)
-a^2\sinc(af_x)\sinc(af_y).
\]
El tercer término resta el cuadrado central contado dos veces.
\textbf{Validación fuerte:} si $a=L$, la cruz \emph{degenera algebraicamente} en un
cuadrado sólido $L\times L$ ---lo verifico numéricamente.

\subsection{Doble cuadrado (Punto 14)}
\textbf{Abertura:} cuadrado de lado $a$ y cuadrado de lado $3a$, separados $2a$
borde-a-borde $\Rightarrow D=4a$ centro-a-centro (geometría fija del enunciado, no
un parámetro libre). Es el bloque 3 con
$A_1=a^2\sinc(af_x)\sinc(af_y)$ y $A_2=9a^2\sinc(3af_x)\sinc(3af_y)$.
\textbf{Validación universal:} $I(0,0)=(\text{área})^2=(a^2+9a^2)^2=(10a^2)^2$.

\subsection{Rendija(s), 1 a $N$ (Puntos 2/6)}
\textbf{Abertura:} $N$ rendijas de ancho $a$, período $d$, idealizadas infinitas en
$y$. El patrón factoriza en \emph{envolvente} $\times$ \emph{factor de red}:
\[
I=\underbrace{\sinc^2\!\Big(\tfrac{a\sin\theta}{\lambda}\Big)}_{\text{una rendija}}
\cdot
\underbrace{\Big[\tfrac{\sin(N\pi d\sin\theta/\lambda)}{N\sin(\pi d\sin\theta/\lambda)}\Big]^2}_{\text{factor de red}} .
\]
\textbf{Validación:} $N=1\Rightarrow$ el factor de red vale 1 (rendija simple);
mínimos de la envolvente en $\sin\theta=m\lambda/a$; con $d=3a$ el orden $m=3$ cae
sobre un cero de la envolvente $\Rightarrow$ \emph{orden ausente}, que el simulador
reproduce.

\subsection{Escalón de Michelson (Punto 15)}
\textbf{Abertura:} escalera de $N$ láminas de vidrio (índice $n$, espesor $h$,
saliente $s$). Entre peldaños vecinos el desfase óptico tiene un aporte geométrico
y uno de vidrio: $\Delta=(n-1)h+s\sin\theta$, y
\[
I=\sinc^2\!\Big(\tfrac{s\sin\theta}{\lambda}\Big)
\Big[\tfrac{\sin(N\pi\Delta/\lambda)}{N\sin(\pi\Delta/\lambda)}\Big]^2 .
\]
\textbf{Qué explico:} el término constante $(n-1)h\sim$ mm empuja los órdenes
principales a $m$ enormes $\Rightarrow$ la altísima dispersión típica del escalón.

\subsection{Dos semicírculos (Punto 9)}
\textbf{Por qué es distinta:} la unión de un semicírculo de radio $r_1$ (arriba) y
otro $r_2$ (abajo) \emph{no} tiene TF cerrada simple. \textbf{Qué hice:} el patrón
2D lo saco por FFT de la máscara, pero el \emph{valor axial} sí es analítico ---en
$x'=y'=0$ toda sinc/Bessel vale 1 y la integral se reduce al área:
\[
\frac{I(0,0)}{I_0}=\Big(\frac{\text{área}}{\lambda z}\Big)^2,\quad
\text{área}=\frac{\pi(r_1^2+r_2^2)}{2}.
\]
\textbf{Validación cruzada:} comparo el centro de la FFT contra esta fórmula
cerrada (el status muestra la diferencia \%). Límites: $r_1=r_2\to$ disco de Airy;
$r_2=0\to$ un semicírculo.

\subsection{Doble círculo --- zonas de Fresnel (Punto 19)}
\textbf{La única figura donde muestro los dos regímenes lado a lado.} Abertura:
disco $r_1$ en tres cuadrantes, extendido a $r_2$ solo en el cuadrante
superior-derecho. Con $r_1=\sqrt{\lambda z}$ (1.\textsuperscript{a} zona) y
$r_2=\sqrt{2\lambda z}$ (2.\textsuperscript{a} zona) calculo Fraunhofer ($|\Fou|^2$)
y Fresnel (\texttt{fresnel\_propagate}) sobre la misma máscara.
\textbf{Validación:} la amplitud axial de Fresnel da $|U|\approx1.5\Rightarrow I=2.25$
(y $I=4$ para una zona completa) ---números que salen de la teoría de zonas.

\subsection{Redes en cascada (Punto 6)}
\textbf{Abertura:} dos redes binarias idénticas ($d=3a$), la segunda desplazada
$s$. \textbf{Idea:} el producto de dos transmisiones binarias
$t_1(x)\,t_2(x-s)$ solo deja pasar donde \emph{ambas} abren $\Rightarrow$ vuelve a
ser una red binaria, y $I=|\Fou\{t_1t_2(\cdot-s)\}|^2$ (FFT 1D).
\textbf{Validación:} $s=0\Rightarrow$ una sola red (pico $=1$); $s=d\Rightarrow$
red de $N-1$ ranuras (factor $((N-1)/N)^2$); $s\in[a,2a]\Rightarrow$ solape nulo,
campo oscuro.

%======================================================================
\section{Cuando no hay fórmula cerrada: la FFT (Código 21)}
%======================================================================

Para formas sin transformada analítica (corazón, hexágono, estrella, triángulo\ldots)
uso el Código 21: rasterizo $t(x,y)$ en una malla y hago $I=|\Fou\{t\}|^2$ con
\texttt{fft2}. \textbf{El punto delicado ---y donde el profesor puede apretar---
es el mapeo de las coordenadas de salida de la FFT a metros.}

Con una ventana de abertura $L_\text{ap}$ muestreada en $N$ puntos:
\begin{equation}
dx=\frac{L_\text{ap}}{N},\qquad df=\frac{1}{L_\text{ap}},\qquad
\boxed{\;dx'=\lambda z\,df=\frac{\lambda z}{L_\text{ap}}\;}
\end{equation}

\begin{decir}
Dos cosas que sé defender aquí: \emph{Primero}, la resolución en la pantalla
depende solo de la ventana física $L_\text{ap}$, no de $N$. Por eso, para que subir
$N$ realmente afine el patrón, hago crecer la ventana con $N$
($L_\text{ap}=N\,dx_\text{ap}$) ---eso es zero-padding, la interpolación estándar de
la FFT. \emph{Segundo}, multiplico la FFT por $dx^2/(\lambda z)$: el $dx^2$ es la
regla de Riemann que convierte la suma discreta de la FFT en la integral continua,
y el $1/\lambda z$ es el prefactor de la ecuación~\eqref{eq:fresnelfft}.
\end{decir}

\subsection{Cómo construyo una abertura sin fórmula: la máscara booleana}

Estas son las figuras ``no esenciales'' (corazón, hexágono, triángulo, estrella):
\textbf{no tienen transformada de Fourier cerrada}, así que el paso 1 de mi receta
---``escribir $t(x,y)$''--- se vuelve literal. En vez de una fórmula, defino la
abertura como un \textbf{predicado booleano}: para cada punto de la malla,
\texttt{True} si está dentro de la figura (transparente) y \texttt{False} si no.
Ese arreglo de 1 y 0 es la abertura, y le aplico exactamente el mismo motor FFT.
Usé tres técnicas según la forma:

\paragraph{(a) Curva implícita --- el corazón.}
Algunas formas se describen por una \emph{desigualdad algebraica}: el interior es
donde una función $F(x,y)\le0$. Para el corazón uso la curva clásica
\[
\Big(\tfrac{x^2}{s^2}+\tfrac{y^2}{s^2}-1\Big)^{\!3}
-\tfrac{x^2}{s^2}\,\tfrac{y^3}{s^3}\;\le\;0 ,
\]
con $s$ el radio característico. Todo punto que cumple la desigualdad es
transparente. En código es una sola línea vectorizada sobre la malla ---sin
bucles--- y por eso el corazón hereda su simetría especular en $x$, que reaparece
en el patrón.

\paragraph{(b) Intersección de semiplanos --- hexágono y triángulo.}
Cualquier \textbf{polígono convexo} es la intersección de semiplanos: un punto está
dentro si y solo si está ``del lado bueno'' de \emph{todos} los lados a la vez. Un
lado con normal $\hat n_k=(\cos\alpha_k,\sin\alpha_k)$ y apotema $r$ define el
semiplano $x\cos\alpha_k+y\sin\alpha_k\le r$; el interior del polígono es el
$\bigwedge_k$ (AND lógico) de todos:
\[
\text{dentro}(x,y)=\bigwedge_{k}\;\big[\,x\cos\alpha_k+y\sin\alpha_k\le r\,\big].
\]
\begin{itemize}[leftmargin=1.4em]
\item \textbf{Hexágono:} basta con 3 direcciones normales ($\alpha_k=0,60,120^\circ$)
usando $|\text{proyección}|\le$ apotema, porque los lados opuestos comparten
dirección; apotema $=s\sqrt3/2$.
\item \textbf{Triángulo equilátero:} 3 semiplanos (normales a los 3 lados),
apotema (inradio) $=s/2$. Como el triángulo \emph{no} tiene centro de inversión,
su patrón adquiere simetría de orden $2\times3=6$ ---la estrella de seis puntas
que menciono abajo.
\end{itemize}
Ambos incluyen un ángulo de giro \texttt{rot}: roto la malla con una matriz de
rotación \emph{antes} de aplicar el test, así giro la figura sin reescribir las
desigualdades.

\paragraph{(c) Radio en función del ángulo --- la estrella.}
Una estrella no es convexa, así que no sirve la intersección de semiplanos. La
defino en \textbf{coordenadas polares}: un punto está dentro si su radio no supera
un radio-frontera $r(\theta)$ que oscila entre el radio exterior $s$ (en cada
punta) y el interior $s\cdot f_\text{int}$ (entre puntas):
\[
\text{dentro}(x,y)=\big[\,r\le r(\theta)\,\big],\qquad
r=\sqrt{x^2+y^2},\;\;\theta=\operatorname{atan2}(y,x),
\]
donde $r(\theta)$ es una onda triangular con \texttt{puntas} periodos en
$[0,2\pi)$. Cambiar el número de picos $p$ es cambiar el periodo de esa onda, y el
patrón responde con simetría $2p$.

\begin{decir}
Lo que quiero que se entienda: para estas figuras \emph{la física no cambia en
nada}. Sigue siendo $I=|\Fou\{t\}|^2$. Lo único distinto es que $t(x,y)$ no viene
de una fórmula de difracción, sino de una condición geométrica ---``¿este píxel
está dentro de la figura?''. Elegí la técnica según la forma: desigualdad
algebraica si la curva es implícita (corazón), AND de semiplanos si es un polígono
convexo (hexágono, triángulo), y comparación de radio si es un contorno polar
(estrella). Después, el mismo FFT para todas.
\end{decir}

\textbf{Las simetrías son mi mejor evidencia de que el mapeo está bien:} un
triángulo equilátero da un patrón de \emph{seis} puntas (clásico y
contraintuitivo), el hexágono da simetría 6, una estrella de $p$ puntas da $2p$.
La TF respeta la geometría, así que si las simetrías salen, el código está bien.

%======================================================================
\section{Fresnel: de dónde sale el método de un solo FFT}
\label{sec:fresnel}
%======================================================================

Aquí está el corazón de lo que el profesor quiere oír sobre Fresnel. \textbf{No
integré numéricamente la doble integral} de la ecuación~\eqref{eq:fresnel}
---eso sería un cuádruple bucle, lentísimo. En vez de eso reconocí que la
ecuación~\eqref{eq:fresnelfft} \emph{ya es una transformada de Fourier}, así que
Fresnel se calcula con \textbf{un solo FFT}:
\begin{equation}
U(x',y')=\frac{1}{i\lambda z}\,
\underbrace{e^{\frac{i\pi}{\lambda z}(x'^2+y'^2)}}_{\text{chirp de salida}}\;
\Fou\Big\{\,t(x,y)\,
\underbrace{e^{\frac{i\pi}{\lambda z}(x^2+y^2)}}_{\text{chirp de entrada}}\Big\}.
\end{equation}

Esto es exactamente la función \texttt{fresnel\_propagate}:
\begin{enumerate}[leftmargin=1.5em]
\item Multiplico la abertura por el \textbf{chirp de entrada} $Q_1=e^{i\pi(X^2+Y^2)/\lambda z}$.
\item Hago \texttt{fft2} (con \texttt{fftshift}/\texttt{ifftshift} para centrar).
\item Multiplico por el \textbf{chirp de salida} $Q_2$ y por el prefactor
$1/(i\lambda z)$, con el paso de salida $dx'=\lambda z/(N\,dx)$.
\end{enumerate}

\begin{decir}
La diferencia entre mi Fresnel y mi Fraunhofer es \emph{una línea}: si dejo el
chirp de entrada $Q_1$, es Fresnel; si lo quito ($Q_1=1$), es Fraunhofer. En el
Código 22 calculo los dos con la misma malla y los superpongo ---por eso puedo
mostrar cómo el patrón de Fresnel converge al de Fraunhofer al alejarme.
\end{decir}

\subsection{El eje angular común: el truco del Código 22}
Para superponer patrones calculados a distintas distancias $z$ (distintos $N_F$)
sin reescalar, uso como abscisa el seno del ángulo
\[
\sin\theta=\frac{x'}{z},
\]
que es \emph{independiente de $z$}. Así la curva de Fraunhofer queda fija y las de
Fresnel convergen a ella cuando $N_F\to0$.

\subsection{Validación contra la espiral de Cornu}
Para una rendija, la integral de Fresnel se expresa con las integrales de Fresnel
(clotoide) $C(v),S(v)$:
\[
I=\frac{I_0}{2}\Big([C(u_2)-C(u_1)]^2+[S(u_2)-S(u_1)]^2\Big),\quad
u_{1,2}=\sqrt{\tfrac{2}{\lambda z}}\Big(\mp\tfrac{a}{2}-x'\Big).
\]
\textbf{Cómo validé:} calculo el patrón con \texttt{fresnel\_propagate} y lo comparo
contra esta fórmula usando \texttt{scipy.special.fresnel}. La correlación es
$\ge0.9999$ a $N_F=2$, y a $N_F=0.1$ el Fresnel y el Fraunhofer son
indistinguibles (correlación $1.0$). Eso demuestra que mi motor de un solo FFT
reproduce la física exacta de Fresnel, no una aproximación mía.

\subsection{Muestreo del chirp (por qué a veces avisa ``submuestreo'')}
El chirp cuadrático oscila cada vez más rápido en los bordes; para no aliasearlo
necesito $N\gtrsim\text{pad}^2N_F$ (Nyquist). La función \texttt{\_pad\_muestreo}
elige la ventana más grande que aún cumpla Nyquist; si tiene que forzar el piso,
el programa lo avisa en rojo. \textbf{Esto lo menciono para mostrar que sé dónde
está el límite numérico del método}, no solo la fórmula.

%======================================================================
\section{El número de Fresnel: cómo sé cuál régimen aplica}
\label{sec:regimen}
%======================================================================

Cada código muestra una \textbf{caja de régimen} que colorea verde o rojo. La lógica:
\begin{equation}
N_F=\frac{D^2}{\lambda z},\qquad
z_\text{min}=\frac{2D^2}{\lambda}\;\Longleftrightarrow\; N_F=0.5,
\end{equation}
con $D$ = mayor extensión lineal de la abertura. El criterio sale de exigir que la
fase del chirp de entrada sea pequeña.

\begin{mdframed}[linecolor=verde,linewidth=1pt,backgroundcolor=grispanel]
\centering
$N_F\le 0.5$: \textcolor{verde}{\textbf{FRAUNHOFER}} (campo lejano) $\Rightarrow$ Códigos 20/21 válidos.\\
$N_F$ grande: \textcolor{rojo}{\textbf{FRESNEL}} (campo cercano) $\Rightarrow$ usar Código 22.
\end{mdframed}

\begin{decir}
Ojo con la intuición al revés que el profesor puede usar para probarme:
$N_F$ \emph{grande} es campo \emph{cercano} (Fresnel), y $N_F$ \emph{pequeño} es
campo \emph{lejano} (Fraunhofer). Puse esa caja de régimen precisamente para no dar
nunca una respuesta físicamente falsa: si estoy en Fresnel, el código me lo dice en
rojo y me remite al Código 22.
\end{decir}

%======================================================================
\section{Validación: cómo sé que no me equivoqué}
\label{sec:validacion}
%======================================================================

Todo el núcleo físico está respaldado por \texttt{validacion\_fisica.py}: una suite
\emph{headless} de \textbf{40 comprobaciones}, cada figura contra un caso límite
conocido. Es mi respuesta a ``¿cómo sabes que está bien?'': no lo sé de vista, lo
sé porque cada expresión pasa su prueba.

\begin{center}
\begin{tabular}{@{}ll@{}}
\toprule
\textbf{Figura} & \textbf{Caso límite que verifico} \\
\midrule
Marco+círculo & $I(0,0)=1$; franjas de paso $\lambda z/D$ \\
Rectángulo & primer cero del sinc en $\lambda z/a$ \\
Círculo & primer anillo de Airy en $1.22\lambda z/2R$ \\
Cruz & degenera en cuadrado si $a=L$ \\
Rect.\ rotado & gira solidario; $\theta=0$ = alineado \\
Doble cuadrado & $I(0,0)=(10a^2)^2$ \\
Red $d=3a$ & orden $m=3$ ausente \\
Zonas de Fresnel & amplitud axial $I=4$ (una zona) \\
Rendija (Fresnel) & correlación $\ge0.9999$ vs.\ Cornu \\
Convergencia & Fresnel$\to$Fraunhofer a $N_F=0.1$ \\
\bottomrule
\end{tabular}
\end{center}
\textbf{Resultado: 40/40 OK.}

%======================================================================
\section{Preguntas que probablemente me hará y cómo respondo}
%======================================================================

\begin{pregunta}
\textcolor{rojo}{\textbf{P:}} \emph{¿Por qué usas \texttt{np.sinc} y no
$\sin(x)/x$?}\\
\textbf{R:} Porque \texttt{np.sinc} es la sinc \emph{normalizada}:
$\sinc(u)=\sin(\pi u)/(\pi u)$, con ceros en enteros. La TF de un rectángulo de
lado $a$ es $a\,\sinc(af)$ sin ningún $\pi$ extra en el argumento. Si usara
$\sin/x$ tendría que meter los $\pi$ a mano y es donde más se equivoca la gente al
traducir fórmulas de un PDF a numpy.
\end{pregunta}

\begin{pregunta}
\textcolor{rojo}{\textbf{P:}} \emph{Tu FFT, ¿te da Fresnel o Fraunhofer?}\\
\textbf{R:} Depende de si conservo el chirp de entrada. La misma \texttt{fft2}:
con la abertura multiplicada por $e^{i\pi(x^2+y^2)/\lambda z}$ da Fresnel; sin ese
factor da Fraunhofer. En el Código 22 calculo ambos con la misma malla justamente
para mostrar esa relación.
\end{pregunta}

\begin{pregunta}
\textcolor{rojo}{\textbf{P:}} \emph{¿Por qué multiplicas la FFT por
$dx^2/(\lambda z)$?}\\
\textbf{R:} Son dos cosas distintas. El $dx^2$ convierte la suma discreta de la FFT
en la integral continua (regla de Riemann: $\sum\to\int dx\,dy$). El $1/(\lambda z)$
es el prefactor $1/(i\lambda z)$ de la integral de difracción ---en intensidad el
$i$ desaparece y queda el módulo.
\end{pregunta}

\begin{pregunta}
\textcolor{rojo}{\textbf{P:}} \emph{¿Por qué el triángulo da seis puntas y no tres?}\\
\textbf{R:} Porque el módulo al cuadrado de la TF impone simetría de inversión:
$|\Fou\{t\}(-k)|=|\Fou\{t\}(k)|$ aunque la figura no la tenga. El triángulo tiene
simetría de orden 3 pero no centro de inversión, así que el patrón adquiere orden
$2\times3=6$. Es la firma de que la FFT respeta la geometría.
\end{pregunta}

\begin{pregunta}
\textcolor{rojo}{\textbf{P:}} \emph{¿La resolución del patrón mejora si subo $N$?}\\
\textbf{R:} No por sí sola: $dx'=\lambda z/L_\text{ap}$ depende de la ventana física,
no de $N$. Por eso hago crecer la ventana con $N$ (zero-padding); así subir $N$ sí
afina el patrón. Es un punto sutil que verifiqué a propósito.
\end{pregunta}

\begin{pregunta}
\textcolor{rojo}{\textbf{P:}} \emph{¿Por qué en el doble cuadrado $D=4a$ y no otro?}\\
\textbf{R:} Es geometría del enunciado, no un parámetro libre. Cuadrados de lado $a$
y $3a$ separados $2a$ borde-a-borde: la distancia centro-a-centro es
$a/2+2a+3a/2=4a$. Ese $D$ fija el período de las franjas de interferencia
($\lambda z/D$).
\end{pregunta}

\begin{pregunta}
\textcolor{rojo}{\textbf{P:}} \emph{Para el marco, ¿por qué restas y no sumas los
dos rectángulos?}\\
\textbf{R:} Porque el marco es el rectángulo exterior \emph{menos} el hueco
interior. Como la TF es lineal, la transformada de la resta de indicadores es la
resta de transformadas: $A_\text{ext}-A_\text{int}$. Es el mismo argumento de
inclusión-exclusión que uso en la cruz.
\end{pregunta}

\begin{pregunta}
\textcolor{rojo}{\textbf{P:}} \emph{¿Cómo validaste el escalón / las zonas de Fresnel
si no tienen un ``primer cero'' obvio?}\\
\textbf{R:} Con el valor axial. En $x'=y'=0$ todas las sinc/Bessel valen 1 y la
integral se reduce al área de la abertura, así que $I(0,0)=(\text{área}/\lambda z)^2$
es un número exacto que puedo comparar contra el centro de mi patrón numérico. Para
las zonas de Fresnel ese número es $I=4$ (una zona) y $2.25$ (la abertura del
punto 19), que salen de la teoría de zonas.
\end{pregunta}

\begin{pregunta}
\textcolor{rojo}{\textbf{P:}} \emph{¿Qué pasa si pongo la abertura muy grande o $z$
muy corto?}\\
\textbf{R:} Caigo en campo cercano: $N_F$ sube por encima de 0.5 y la caja de
régimen se pone roja avisando que el patrón de Fraunhofer ya no es válido y que
debo usar el Código 22. El código sigue dibujando, pero me advierte que la física
mostrada no corresponde a ese régimen.
\end{pregunta}

\vspace{4pt}
\begin{decir}
\textbf{Cierre.} Si tengo que resumir la defensa: no repetí una deducción distinta
por figura. Implementé \emph{una sola} maquinaria ---transformada de Fourier de la
abertura, con o sin chirp según el régimen--- y para cada figura solo cambié la
$t(x,y)$ que le entra. Cada expresión tiene un caso límite que la valida, y los 40
checks confirman que las 9+9+6 figuras salen correctas.
\end{decir}

\end{document}
